
%\needspace{4\baselineskip}
\color{uniblue}{\subsection{Защита минимального напряжения (ЗМН)}}
\color{black}

\begin{enumerate}

\item
Защита минимального напряжения необходима для выявления снижения напряжения ниже допустимых пределов. В устройстве защита минимального напряжения представлена функциональным блоком «ЗМН», состоящим из независимых функций (ступеней) «ЗМН 1» и «ЗМН 2». ИО ступеней реагируют на симметричное снижение трех линейных напряжений ниже уставки.

\item
Сигнал пуска функции осуществляется при одновременном наличии следующих условий:
\begin{itemize}
    \item функция ЗМН введена в работу;
    \item снижение всех трех линейных напряжений ниже уставки;
    \item отсутствие сигнала блокировки ЗМН из логики БНН;
    \item наличие входного дискретного сигнала разрешения ЗМН (для реализации логики АВР), если введен контроль наличия внешнего разрешения ЗМН уставкой SGF2.    
\end{itemize}

\item
При условии, что ключ «ЗМН1(2)ст. на сигн» введен, действие функции переводится на сигнал, при этом выходной сигнал «ЗМН1(2)ст: Срабатывание» блокируется.

\item
При наличии пусковых условий, срабатывание функции происходит после отсчета выдержки времени таймера T1.

\item
Ввод 1,2 ступени ЗМН в работу на этапе параметрирования устройства осуществляется своим программным переключателем SGF1 «Ввод\_функции». Информация о введенной ступени отображается выходным сигналом «ЗМН 1(2)ст: Ввод».

\item
Ступени ЗМН могут быть выведены из работы оперативно путем активации общего сигнала «ОВ ЗМН», либо по отдельности каждая активацией сигнала «ОВ ЗМН 1ст» или «ОВ ЗМН 2ст». Выведенное состояние ступени характеризуется активным состоянием сигнала «ЗМН 1: Оперативный вывод», «ЗМН 2: Оперативный вывод» для первой и второй ступени ЗМН соответственно.

\item
При активном состоянии входного сигнала «ЗМН 1 на сигн», «ЗМН 2 на сигн» действие соответствующей ступени ЗМН переводится на сигнал.

\item
Алгоритм работы ЗМН выполнен в соответствии с рисунком \ref{zmn:img1}. В таблице \ref{zmn:tbl1} приведены параметры, необходимые для настройки ЗМН.

\vspace{3mm}
\begin{figure}[H]
\centering
\includegraphics[width=1\textwidth,height=1\textheight,keepaspectratio]{\fbpath/201. ЗМН ТН/_latex/img/zmn.pdf}
\captionof{figure}{Функциональная схема ступени ЗМН}\label{zmn:img1}
\end{figure}


\end{enumerate}
